\section{Introduction}
\label{sec:intro}

% --- Para 1: VLAs are impressive but fragile under distribution shift ---
Vision-language-action models (VLAs) that generate actions via flow matching~\cite{black2024pi0,shukor2025smolvla,lipman2023flow} have reached high success rates on in-distribution benchmarks such as LIBERO~\cite{liu2023libero}.
Yet this near-saturation masks a fragile reality: visual and physical perturbations routinely encountered in deployment---camera viewpoint shifts, object layout changes, robot initial state variation---can degrade success rates by over 65 percentage points---for example, $\pi_0$ drops from 94.2\% to 6.6\% under robot initial-state perturbation, and UniVLA from 95.2\% to 4.3\% under camera viewpoint changes~\cite{fei2025liberoplus}.
Closing this robustness gap without retraining the base policy is a prerequisite for deploying VLAs outside tightly controlled settings.

% --- Para 2: Existing approaches and the gap ---
Two lines of work address VLA robustness.
\emph{Training-time methods} improve the base policy itself through robustness-aware objectives~\cite{srpo2025,robustvla2025,digflow2025} or test-time reinforcement learning~\cite{liu2026ttvla}.
These approaches are effective but require retraining or per-task RL, and their gains are bounded by the training distribution.
\emph{Inference-time methods} leave the policy frozen.
ReconVLA~\cite{chen2026reconvla} applies conformal quantile regression to calibrate action-level uncertainty bounds and detects out-of-distribution states, taking a complementary UQ-driven approach rather than learning explicit corrective residuals.
VLA-Pilot~\cite{li2025vlapilot} applies evolutionary optimization over action samples, but incurs 2.41\,s per control step---too slow for real-time manipulation.
A2C2~\cite{a2c2_2025} learns a per-step correction head for pre-generated action chunks on a frozen VLA, but applies correction unconditionally without multi-sample selection or explicit magnitude control.
In short, no existing inference-time method combines multi-sample conformal selection with bounded residual correction in a two-stage pipeline.

% --- Para 3: Our approach — sample-then-correct ---
We introduce \textbf{ACR-VLA} (\textbf{A}daptive \textbf{C}onformal \textbf{R}efinement for VLA), an inference-time module that improves frozen VLA actions through a \emph{sample-then-correct} pipeline.
At each control step, ACR-VLA draws $K{=}10$ action samples from the frozen VLA, selects the most conforming candidate via Mahalanobis-based nonconformity scoring with shrinkage covariance estimation~\cite{ledoit2004well,braun2025cgs}, and then applies a learned residual correction via a lightweight \emph{Conformal Residual Module} (CRM).
CRM is a 330K-parameter MLP conditioned on the VLA's frozen observation embedding.
It predicts a small residual $\boldsymbol{\delta}$ whose $L_2$ norm is clipped to a bound $\varepsilon$, producing the corrected action $\mathbf{a}^* = a_{\mathrm{base}} + \mathrm{clip}(\boldsymbol{\delta})$.
The critical design insight is that correction operates on a high-quality starting point---the conformal-selected action---and that bounding the correction magnitude prevents compounding errors from destabilizing the closed-loop trajectory.
The base VLA is never modified; ACR-VLA attaches as a post-hoc module to any frozen flow-matching VLA.

% --- Para 4: Key findings (with exp_id references) ---
We evaluate ACR-VLA on LIBERO-Object (10 tasks, 100 trials per method) and LIBERO-Spatial (10 tasks, 100 trials) using SmolVLA~\cite{shukor2025smolvla} as the frozen backbone. Three findings emerge:

\textbf{(1)~Sample-then-correct outperforms sample-then-select.}
On LIBERO-Object, observation-conditioned CRM achieves 19.0\% success versus 12.0\% for conformal selection alone---a +7.0 percentage point improvement representing a 58\% relative gain.
% exp_id: eval-acr-object-v4-obs-only (19.0%), eval-argmin-only-object (12.0%)
% TODO: update with pooled seed=42+43 results and Fisher exact p-value
Averaging all $K$ samples (VLA-mean) or using a single sample (VLA-only) both yield 12.0\%, confirming that the gain comes from learned correction, not from increased sampling alone.
% exp_id: eval-vla-only-object-k1 (12.0%), eval-vla-mean-object (12.0%)

\textbf{(2)~Observation features suffice for correction; conformal features do not help.}
Counter to intuition from synthetic experiments, conditioning CRM on per-DOF conformal statistics (perdof: 11.0\%) performs no better than selection alone (12.0\%), while scalar conformal conditioning (9.0\%) is actively harmful.
% exp_id: eval-acr-object-v4-small (perdof 11.0%), eval-acr-object-v4-scalar (9.0%)
Only observation-conditioned CRM (19.0\%) produces meaningful improvement.
This honest negative result highlights a gap between per-step prediction accuracy---where richer features help in offline evaluation---and closed-loop task success, where covariance estimation noise from $K{=}10$ samples undermines per-DOF conditioning.

\textbf{(3)~Correction magnitude control is critical and previously unrecognized.}
Without magnitude bounding, the per-DOF-conditioned CRM's natural output ($\|\boldsymbol{\delta}\|_2 \approx 1.42$) causes catastrophic failure: 0\% success on LIBERO-Spatial, as compounding correction errors rapidly drive the state distribution away from the CRM's training regime.
Constraining corrections to $\|\boldsymbol{\delta}\|_2 \leq 0.1$ recovers the full +7.0\,pp gain.
% exp_id: eval-acr-spatial-v4-default-r3 (norm=1.42, 0%), eval-acr-object-v4-obs-only (norm=0.1, 19.0%)
Observation-conditioned CRM maintains stable 35--36\% success on LIBERO-Spatial across $\varepsilon \in [0.05, 0.3]$, while per-DOF-conditioned CRM collapses from 35\% to 26\% over the same range---further evidence that observation features provide more robust conditioning for learned correction.
% exp_id: eval-mag-obsonly-norm005, eval-mag-obsonly-norm03, eval-mag-perdof-norm03

We note that correction offers targeted improvements on object-centric manipulation rather than universal gains: on LIBERO-Spatial, the best CRM variant performs within 1\,pp of the selection baseline, suggesting that the correction mechanism is most effective when distribution shift induces systematic action biases that a learned module can identify and compensate.

% --- Para 5: Contributions ---
Our contributions are:
\begin{enumerate}[leftmargin=*,nosep]
  \item \textbf{A sample-then-correct pipeline for frozen VLAs.}
  We propose ACR-VLA, which combines conformal sample selection with learned residual correction for frozen flow-matching VLAs.
  On LIBERO-Object, this pipeline improves SmolVLA success by +7.0\,pp over selection alone.
  % TODO: update with pooled seed results and p-value: (p < 0.05 on combined 200 trials, Fisher exact test)

  \item \textbf{Systematic conditioning comparison and honest negative result.}
  We compare three CRM conditioning variants---per-DOF conformal features, scalar uncertainty, and observation features---and find that only observation conditioning is effective.
  This challenges the intuition that richer uncertainty features should yield better corrections, and reveals that covariance estimation noise at $K{=}10$ undermines per-DOF conditioning in closed-loop settings.

  \item \textbf{Systematic empirical study of correction magnitude sensitivity for VLA action refinement.}
  We show that unbounded correction causes catastrophic failure and that observation-conditioned CRM is uniquely robust across a $6{\times}$ range of magnitude bounds ($\varepsilon \in [0.05, 0.3]$), establishing magnitude control as a critical design axis for inference-time action correction.
\end{enumerate}
